\section{引言}
\label{sec:method}
编译系统负责将高级语言编写的源程序转换为目标平台能够执行的程序，其处理过程涉及预处理、编译、汇编和链接等多个阶段。其中，编译器内部又包含词法分析、语法分析、语义分析、中间代码生成、代码优化和目标代码生成等环节。不同阶段承担着不同的任务，其输出也具有不同的表示形式。因此，要理解编译系统的完整工作机制，不仅需要掌握各阶段的基本功能，还需要结合实际编译产物，分析源程序如何逐步转换为可执行程序。

本实验以 GCC 和 Clang/LLVM 为主要研究对象，首先选取包含宏定义、条件分支、循环及函数调用的 C 语言程序，利用编译器命令行选项分别获取预处理文件、中间表示、汇编文件和目标文件，分析各阶段的具体功能及其输入输出关系。同时，通过比较不同优化选项和调试选项下的编译结果，进一步研究编译优化对控制流、中间表示和目标代码结构的影响。

在此基础上，本实验设计了两个具有不同语言特征的 SysY 示例程序，涵盖整数与浮点运算、变量赋值、条件分支、循环、函数调用、数组及作用域等语言特性。针对相同的源程序语义，分别编写 LLVM IR 和 RISC-V 汇编实现，分析控制流、静态单赋值形式（SSA）、函数调用约定及数据访问等机制，并通过链接 SysY 运行时库、生成目标程序和执行测试，检验不同层次程序表示之间的语义一致性。

此外，为进一步了解面向人工智能加速器的编译技术，本实验以 AscendNPU IR 官方 VecAdd 示例为对象，探索基于 MLIR 的多级中间表示及逐层降低（lowering）过程，重点分析核类型推导、片上内存规划、跨流水同步和设备操作模板转换等中端处理机制。

本实验通过编译产物对照、程序运行验证与中间表示分析相结合的方法，研究传统编译工具链、LLVM IR 与目标汇编之间的关系，并进一步认识面向专用硬件的多级中间表示转换过程。
